% !TEX program = xelatex
% ============================================================
%  YOLOv8 + OpenVINO 部署项目 —— 简历条目（可直接编译）
%  编译：xelatex 简历-项目.tex
%  用法：
%   1) 只想要内容：复制 \section* 到 \end{document} 之间的部分，
%      粘进你现有的简历模板即可（不依赖本文档的导言区）。
%   2) 想自己看效果：直接编译本文件。
%   3) 版面紧张时用文末「精简版」或「一行版」。
% ============================================================
\documentclass[11pt]{ctexart}

\usepackage[a4paper,margin=2.2cm]{geometry}
\usepackage{enumitem}
\usepackage{xcolor}
\usepackage{hyperref}
\hypersetup{colorlinks=true, linkcolor=blue!45!black, urlcolor=blue!45!black}

\setlist[itemize]{leftmargin=1.3em, itemsep=2pt, topsep=2pt, parsep=0pt}
\pagestyle{empty}
\setlength{\parindent}{0pt}

% 项目标题 + 时间
\newcommand{\projecttitle}[2]{%
  \textbf{#1}\hfill{\small #2}\\[2pt]%
}

\begin{document}

\section*{项目经历}

\projecttitle{YOLOv8 + OpenVINO 端到端部署（CPU 推理与 INT8 量化）}{2026.10}
{\small 个人项目 \quad 代码：\href{https://github.com/liumaker/yolov8-openvino-deploy}{github.com/liumaker/yolov8-openvino-deploy}}

\begin{itemize}
  \item \textbf{导出与转换}：打通 PyTorch\,$\rightarrow$\,ONNX（opset 12，静态/动态轴）\,$\rightarrow$\,OpenVINO IR（fp32/fp16）链路，
        用 Netron 核对计算图；以同一输入张量对齐三端原始输出，相对误差 $3.2\times10^{-6}$，证明导出无语义损失。

  \item \textbf{自研前后处理}：不依赖 Ultralytics 高层 API，手写 letterbox、输出解码（84$\,=\,$4$\,+\,$80，anchor-free）与按类 NMS，
        与官方实现逐框对齐（5/5 框匹配，坐标偏差 $0.0001$\,px），单帧端到端 $56.9$\,ms。

  \item \textbf{INT8 量化}：用 NNCF 做 PTQ（300 张校准图）。朴素量化掉点严重（$-5.9$ AP），改用 accuracy-aware 逐层回退，
        COCO val2017 500 张上 mAP $0.3998\rightarrow0.3886$（$-1.1$ 点，$-2.8\%$），
        单帧延迟 $30.7\rightarrow16.6$\,ms（$-46\%$），吞吐 $32\rightarrow57$\,FPS。

  \item \textbf{并发与工程化}：benchmark\_app 对比 FP32/FP16/INT8 延迟与吞吐；实现多路视频流并发（每路独立 \texttt{InferRequest}），
        4 路 INT8 总吞吐 $48$\,FPS（约为 FP32 的 $1.9\times$），进程 CPU $\sim630\%$，定位瓶颈为 CPU 算力；
        沉淀 10 个可复现脚本与自动化验收（数值对齐 / mAP / 基准）。
\end{itemize}

{\small \textbf{技术栈}：Python · PyTorch · ONNX / ONNX Runtime · OpenVINO · NNCF · OpenCV · pycocotools}

% \vspace{0.6em}
% \rule{\linewidth}{0.4pt}

% ============================================================
%  精简版（版面紧张时替换上面的 itemize）
% ============================================================
% \begin{itemize}
%   \item 完成 YOLOv8n 从 PyTorch 到 OpenVINO INT8 的 CPU 部署：ONNX 导出（opset 12、动态轴）、IR 转换、
%         自研 letterbox/解码/NMS，与官方输出逐框一致（坐标误差 $<0.001$\,px）。
%   \item NNCF PTQ（300 张校准图）+ accuracy-aware 逐层回退：mAP $0.3998\rightarrow0.3886$（$-1.1$ 点），
%         单帧延迟 $30.7\rightarrow16.6$\,ms（$-46\%$）；4 路视频流并发 $48$\,FPS（FP32 $25$\,FPS）。
% \end{itemize}
%
% ============================================================
%  一行版（放在技能栏或项目列表的一行）
% ============================================================
% YOLOv8 + OpenVINO CPU 部署：INT8 量化后单帧延迟 $30.7\rightarrow16.6$\,ms（$-46\%$），
% mAP 仅降 $1.1$ 点；4 路视频流并发 $48$\,FPS。

\end{document}
